\documentclass[11pt]{article}
\usepackage{fancyhdr}
\usepackage[table]{xcolor}
\usepackage[margin=0.7in]{geometry}
\usepackage{caption}
\usepackage{hyperref}
\usepackage{graphicx}
\usepackage{float}
\usepackage{booktabs}
\usepackage{colortbl}

% Set header and footer
\pagestyle{fancy}
\fancyhf{}
\fancyhead[L]{USC Facilities: Rain Barrel Water Pump System[Handout 1]}
\fancyfoot[C]{\thepage}

\begin{document}

% Table 1a
\begin{table}[H]
\centering
\caption*{Table 1a: Operational efficiency needs.}
\label{table:1a}
\rowcolors{2}{lightgray}{}
\begin{tabular}{|p{12cm}|c|c|c|c|c|}
\hline
\textbf{Need} & \textbf{\#1} & \textbf{\#2} & \textbf{\#3} & \textbf{\#4} & \textbf{Sum} \\
\hline
1.1. The system should minimize the time required to transfer water from the tank to the garden. & 5 & 6 & 7 & 4 & \cellcolor{yellow}22 \\
1.2. The system must operate effectively with minimal manual intervention. & 0 & 1 & 0 & 0 & 1 \\
1.3. Water distribution should cover the entire garden area efficiently. & 4 & 4 & 2 & 5 & \cellcolor{yellow}15 \\
1.4. The system should be able to operate without frequent start-stop interruptions. & 0 & 0 & 0 & 1 & 1 \\
1.5. The system should ensure minimal water wastage during operation. & 3 & 1 & 0 & 2 & 6 \\
1.6. Water output should be consistent and reliable during the watering process. & 3 & 3 & 2 & 3 & \cellcolor{yellow}11 \\
1.7. The system should be flexible enough to handle varying watering schedules. & 0 & 0 & 4 & 1 & 5 \\
1.8. The system should be scalable to handle garden expansion or additional tanks. & 0 & 1 & 2 & 2 & 5 \\
1.9. The system should allow for easy monitoring of water usage and distribution. & 0 & 0 & 0 & 1 & 1 \\
1.10. The system should be able to adapt to water demand and supply conditions. & 0 & 0 & 0 & 0 & 0 \\
1.11. The system should maintain functionality during peak water demand periods. & 2 & 1 & 2 & 1 & 6 \\
1.12. The system should be easy to activate and deactivate when needed. & 3 & 2 & 5 & 2 & \cellcolor{yellow}12 \\
1.13. The system should be adaptable to changes in garden layout or design. & 0 & 0 & 0 & 0 & 0 \\
1.14. The system should be able to handle variable water volumes effectively. & 1 & 1 & 0 & 2 & 4 \\
1.15. The system should minimize the physical effort required to operate it. & 3 & 2 & 2 & 1 & \cellcolor{yellow}8 \\
1.16. The system should be easy to set up and configure for different garden zones. & 0 & 0 & 0 & 1 & 1 \\
1.17. The system should have the capability to adjust flow according to plant needs. & 0 & 0 & 0 & 0 & 0 \\
1.18. The system should be capable of rapid deployment or removal as needed. & 0 & 0 & 0 & 0 & 0 \\
1.19. The system should facilitate quick troubleshooting during operational issues. & 3 & 1 & 5 & 2 & \cellcolor{yellow}11 \\
1.20. The system should be able to resume operations seamlessly. & 2 & 1 & 1 & 1 & 5 \\
\hline
\end{tabular}
\end{table}

% Table 1b
\begin{table}[H]
\centering
\caption*{Table 1b: Pressure and flow needs.}
\label{table:1b}
\rowcolors{2}{lightgray}{}
\begin{tabular}{|p{12cm}|c|c|c|c|c|}
\hline
\textbf{Need} & \textbf{\#1} & \textbf{\#2} & \textbf{\#3} & \textbf{\#4} & \textbf{Sum} \\
\hline
2.1. The system should provide sufficient pressure for effective water delivery. & 3 & 2 & 5 & 2 & \cellcolor{yellow}12 \\
2.2. The system should maintain consistent pressure regardless of elevation. & 2 & 1 & 0 & 1 & 4 \\
2.3. The system should prevent significant pressure drops during operation. & 1 & 0 & 0 & 0 & 1 \\
2.4. The system should allow adjustment of pressure for different applications. & 0 & 0 & 0 & 0 & 0 \\
2.5. The system should provide feedback on water pressure for user awareness. & 0 & 0 & 0 & 0 & 0 \\
2.6. The system should be capable of maintaining pressure over long distances. & 2 & 2 & 0 & 3 & 7 \\
2.7. The system should ensure even water distribution across all garden areas. & 0 & 0 & 2 & 1 & 3 \\
2.8. The system should avoid pressure surges that could damage components. & 1 & 1 & 3 & 1 & 6 \\
2.9. The system should function effectively under varying pressure conditions. & 1 & 1 & 0 & 0 & 2 \\
2.10. The system should provide adequate flow rate for all intended watering methods. & 3 & 3 & 2 & 0 & \cellcolor{yellow}8 \\
2.11. The system should allow easy detection and correction of flow issues. & 0 & 0 & 0 & 0 & 0 \\
2.12. The system should be capable of maintaining flow in low-pressure scenarios. & 0 & 1 & 6 & 0 & 7 \\
2.13. The system should operate efficiently with varying water levels in the tank. & 4 & 4 & 2 & 3 & \cellcolor{yellow}13 \\
2.14. Water should reach all parts of the garden without pressure loss. & 3 & 3 & 0 & 2 & \cellcolor{yellow}8 \\
2.15. The system should support gradual changes in pressure to avoid spikes. & 1 & 1 & 0 & 0 & 2 \\
2.16. The system should allow for easy verification of flow performance. & 0 & 0 & 0 & 0 & 0 \\
2.17. The system should ensure minimal disruption to flow during adjustments. & 0 & 0 & 0 & 0 & 0 \\
2.18. The system should be able to function at both high and low flow rates. & 0 & 0 & 0 & 0 & 0 \\
2.19. The system should facilitate consistent flow under fluctuating conditions. & 1 & 2 & 0 & 2 & 5 \\
2.20. The system should maintain optimal flow for the intended irrigation method. & 3 & 3 & 0 & 0 & 6 \\
\hline
\end{tabular}
\end{table}

% Table 1c
\begin{table}[H]
\centering
\caption*{Table 1c: Usability needs.}
\label{table:1c}
\rowcolors{2}{lightgray}{}
\begin{tabular}{|p{12cm}|c|c|c|c|c|}
\hline
\textbf{Need} & \textbf{\#1} & \textbf{\#2} & \textbf{\#3} & \textbf{\#4} & \textbf{Sum} \\
\hline
3.1. The system should be easy for non-technical users to operate. & 4 & 3 & 6 & 4 & \cellcolor{yellow}17 \\
3.2. The system should provide clear instructions for use. & 3 & 2 & 1 & 3 & \cellcolor{yellow}9 \\
3.3. The system should be intuitive to set up and use with minimal training. & 2 & 2 & 2 & 1 & 7 \\
3.4. The system should be accessible to users with different physical capabilities. & 0 & 0 & 0 & 0 & 0 \\
3.5. The system should provide clear feedback on its operational status. & 0 & 0 & 0 & 0 & 0 \\
3.6. The system should be easy to start, stop, and adjust as needed. & 4 & 5 & 5 & 2 & \cellcolor{yellow}16 \\
3.7. The system should allow for quick learning and familiarization. & 2 & 2 & 2 & 1 & 7 \\
3.8. The system should be straightforward to configure for different watering needs. & 2 & 2 & 0 & 0 & 4 \\
3.9. The system should provide visible indicators of key operational parameters. & 0 & 0 & 0 & 0 & 0 \\
3.10. The system should minimize the number of steps required for operation. & 3 & 3 & 1 & 4 & \cellcolor{yellow}11 \\
3.11. The system should have easy-to-read and understand operational guidelines. & 1 & 1 & 0 & 1 & 3 \\
3.12. The system should be user-friendly in both daylight and low-light conditions. & 0 & 1 & 0 & 1 & 2 \\
3.13. The system should have clear labeling for all user-accessible components. & 1 & 0 & 0 & 0 & 1 \\
3.14. The system should allow for easy monitoring of its performance and settings. & 0 & 0 & 0 & 0 & 0 \\
3.15. The system should require minimal physical exertion to operate. & 2 & 1 & 0 & 1 & 4 \\
3.16. The system should be easy to reset or reconfigure if needed. & 0 & 0 & 1 & 0 & 1 \\
3.17. The system should facilitate quick and easy shutdown in case of emergency. & 0 & 1 & 0 & 1 & 2 \\
3.18. The system should be easy to use without the need for tools or equipment. & 2 & 1 & 1 & 3 & 7 \\
3.19. The system should allow users to operate it comfortably for extended periods. & 1 & 1 & 0 & 1 & 3 \\
3.20. The system should be designed to minimize errors during operation. & 2 & 2 & 4 & 2 & \cellcolor{yellow}10 \\
\hline
\end{tabular}
\end{table}

% Table 1d
\begin{table}[H]
\centering
\caption*{Table 1d: Maintenance needs.}
\label{table:1d}
\rowcolors{2}{lightgray}{}
\begin{tabular}{|p{12cm}|c|c|c|c|c|}
\hline
\textbf{Need} & \textbf{\#1} & \textbf{\#2} & \textbf{\#3} & \textbf{\#4} & \textbf{Sum} \\
\hline
4.1. The system should require minimal maintenance for reliable operation. & 4 & 5 & 3 & 2 & \cellcolor{yellow}14 \\
4.2. The system should allow easy access to components that need maintenance. & 0 & 1 & 0 & 1 & 2 \\
4.3. The system should have a low frequency of maintenance interventions. & 3 & 2 & 0 & 1 & 6 \\
4.4. The system should allow for identification of components that require servicing. & 0 & 0 & 0 & 1 & 1 \\
4.5. The system should be easy to clean and maintain regularly. & 2 & 1 & 0 & 1 & 4 \\
4.6. The system should provide indications when maintenance is needed. & 0 & 0 & 0 & 0 & 0 \\
4.7. The system should be designed to minimize the likelihood of component failure. & 1 & 2 & 2 & 2 & 7 \\
4.8. The system should be robust against wear and tear in normal conditions. & 3 & 2 & 3 & 1 & \cellcolor{yellow}9 \\
4.9. The system should allow for easy replacement of consumable parts. & 1 & 0 & 0 & 0 & 1 \\
4.10. The system should have a simple process for performing routine maintenance. & 0 & 0 & 0 & 1 & 1 \\
4.11. The system should be able to continue functioning with minor faults or wear. & 1 & 1 & 4 & 0 & 6 \\
4.12. The system should include features that prevent issues from escalating. & 0 & 0 & 0 & 0 & 0 \\
4.13. The system should be resistant to clogging and other common issues. & 1 & 2 & 0 & 0 & 3 \\
4.14. The system should minimize the need for specialized maintenance or skills. & 0 & 0 & 1 & 1 & 2 \\
4.15. The system should be durable under both indoor and outdoor conditions. & 0 & 0 & 0 & 0 & 0 \\
4.16. The system should have a clear and accessible maintenance log. & 0 & 0 & 0 & 1 & 1 \\
4.17. The system should allow for easy troubleshooting of common issues. & 0 & 1 & 0 & 1 & 2 \\
4.18. The system should be easy to drain and store if needed. & 0 & 0 & 0 & 0 & 0 \\
4.19. The system should have protection against damage from external elements. & 0 & 0 & 7 & 0 & 7 \\
4.20. The system should be easy to prepare for seasonal changes in weather. & 0 & 1 & 0 & 1 & 2 \\
\hline
\end{tabular}
\end{table}

% Table 1e
\begin{table}[H]
\centering
\caption*{Table 1e: Environmental impact needs.}
\label{table:1e}
\rowcolors{2}{lightgray}{}
\begin{tabular}{|p{12cm}|c|c|c|c|c|}
\hline
\textbf{Need} & \textbf{\#1} & \textbf{\#2} & \textbf{\#3} & \textbf{\#4} & \textbf{Sum} \\
\hline
5.1. The system should minimize its impact on the local environment. & 2 & 1 & 1 & 1 & 5 \\
5.2. The system should prioritize the use of sustainable and recyclable materials. & 1 & 1 & 2 & 1 & 5 \\
5.3. The system should be energy-efficient and use minimal power. & 2 & 2 & 0 & 2 & 6 \\
5.4. The system should not introduce contaminants to the water or soil. & 3 & 2 & 0 & 2 & 7 \\
5.5. The system should be designed to reduce water wastage. & 0 & 1 & 0 & 2 & 3 \\
5.6. The system should align with sustainability goals of the project site. & 2 & 2 & 2 & 0 & 6 \\
5.7. The system should not disrupt local wildlife or their habitats. & 0 & 0 & 0 & 1 & 1 \\
5.8. The system should be compatible with existing eco-friendly practices on-site. & 1 & 1 & 2 & 0 & 4 \\
5.9. The system should support the collection and use of renewable resources. & 0 & 1 & 1 & 1 & 3 \\
5.10. The system should be designed to minimize noise pollution. & 3 & 2 & 4 & 2 & \cellcolor{yellow}11 \\
5.11. The system should be easily integrated into the existing infrastructure. & 0 & 0 & 0 & 0 & 0 \\
5.12. The system should be resistant to environmental factors like extreme weather. & 0 & 1 & 0 & 0 & 1 \\
5.13. The system should be designed to minimize its carbon footprint. & 2 & 2 & 0 & 2 & 6 \\
5.14. The system should promote the conservation of natural resources. & 0 & 0 & 0 & 0 & 0 \\
5.15. The system should support practices that reduce soil erosion. & 0 & 0 & 0 & 0 & 0 \\
5.16. The system should be able to adapt to changing environmental regulations. & 0 & 0 & 0 & 0 & 0 \\
5.17. The system should not contribute to the pollution of local water sources. & 0 & 1 & 0 & 2 & 3 \\
5.18. The system should use materials that are safe for human interaction. & 1 & 2 & 3 & 2 & \cellcolor{yellow}8 \\
5.19. The system should be compatible with rainwater harvesting practices. & 3 & 4 & 0 & 1 & \cellcolor{yellow}8 \\
5.20. The system should last many years with minimal environmental impact. & 2 & 2 & 5 & 1 & \cellcolor{yellow}10 \\
\hline
\end{tabular}
\end{table}

% Table 1f
\begin{table}[H]
\centering
\caption*{Table 1f: Cost and budgetary needs.}
\label{table:1f}
\rowcolors{2}{lightgray}{}
\begin{tabular}{|p{12cm}|c|c|c|c|c|}
\hline
\textbf{Need} & \textbf{\#1} & \textbf{\#2} & \textbf{\#3} & \textbf{\#4} & \textbf{Sum} \\
\hline
6.1. The system should be cost-effective to implement within the given budget. & 2 & 1 & 4 & 1 & \cellcolor{yellow}8 \\
6.2. The system should have low operational costs once installed. & 1 & 1 & 0 & 1 & 3 \\
6.3. The system should not exceed the specified maximum budget. & 2 & 1 & 1 & 2 & 6 \\
6.4. The system should prioritize cost-saving without compromising performance. & 0 & 0 & 0 & 1 & 1 \\
6.5. The system should allow for flexible budgeting and cost adjustments. & 0 & 0 & 0 & 0 & 0 \\
6.6. The system should avoid ongoing costs that are disproportionately high. & 0 & 0 & 0 & 1 & 1 \\
6.7. The system should leverage available resources to minimize expenses. & 0 & 0 & 0 & 0 & 0 \\
6.8. The system should provide value for money in terms of durability and efficiency. & 1 & 1 & 2 & 1 & 5 \\
6.9. The system should minimize the need for expensive consumables or parts. & 1 & 1 & 2 & 1 & 5 \\
6.10. The system should have a predictable cost structure for maintenance. & 0 & 0 & 3 & 0 & 3 \\
6.11. The system should facilitate potential cost recovery through operation. & 0 & 0 & 0 & 0 & 0 \\
6.12. The system should include a breakdown of anticipated costs and benefits. & 0 & 1 & 0 & 1 & 2 \\
6.13. The system should allow for potential upgrades within the existing budget. & 0 & 0 & 0 & 0 & 0 \\
6.14. The system should be designed with cost-effectiveness in mind from the start. & 0 & 0 & 2 & 0 & 2 \\
6.15. The system should be easy to scale without significant cost increases. & 0 & 0 & 1 & 0 & 1 \\
6.16. The system should not incur unexpected costs during installation or operation. & 0 & 0 & 0 & 0 & 0 \\
6.17. The system should utilize durable components to reduce replacement costs. & 0 & 0 & 4 & 0 & 4 \\
6.18. The system should be designed to minimize financial risk to stakeholders. & 0 & 0 & 0 & 0 & 0 \\
6.19. The system should include provisions for cost-efficient disposal or recycling. & 0 & 0 & 0 & 0 & 0 \\
6.20. The system should consider long-term financial implications of its use. & 0 & 1 & 0 & 1 & 2 \\
\hline
\end{tabular}
\end{table}

% Table 1g
\begin{table}[H]
\centering
\caption*{Table 1g: Safety and compliance needs.}
\label{table:1g}
\rowcolors{2}{lightgray}{}
\begin{tabular}{|p{12cm}|c|c|c|c|c|}
\hline
\textbf{Need} & \textbf{\#1} & \textbf{\#2} & \textbf{\#3} & \textbf{\#4} & \textbf{Sum} \\
\hline
7.1. The system should comply with all safety standards. & 3 & 2 & 2 & 2 & \cellcolor{yellow}9 \\
7.2. The system should be safe to use in various weather. & 0 & 1 & 0 & 2 & 3 \\
7.3. The system should include measures to prevent accidental misuse. & 2 & 3 & 0 & 2 & 7 \\
7.4. The system should be free from sharp edges or hazardous components. & 1 & 1 & 2 & 1 & 5 \\
7.5. The system should not pose any risks to pedestrians or garden users. & 1 & 2 & 2 & 3 & \cellcolor{yellow}8 \\
7.6. The system should have clear labeling to indicate potential hazards. & 1 & 1 & 0 & 2 & 4 \\
7.7. The system should prevent water contamination that could pose health risks. & 0 & 0 & 0 & 0 & 0 \\
7.8. The system should be able to withstand accidental impacts without failure. & 0 & 0 & 0 & 1 & 1 \\
7.9. The system should not obstruct emergency access or exits. & 0 & 0 & 0 & 1 & 1 \\
7.10. The system should have measures to prevent unauthorized access. & 0 & 0 & 0 & 1 & 1 \\
7.11. The system should comply with local environmental regulations. & 0 & 0 & 0 & 1 & 1 \\
7.12. The system should not interfere with existing infrastructure or utilities. & 0 & 0 & 0 & 1 & 1 \\
7.13. The system should be designed to prevent tripping or slipping hazards. & 1 & 1 & 0 & 1 & 3 \\
7.14. The system should be stable and secure when in use and when not in use. & 2 & 1 & 0 & 1 & 4 \\
7.15. The system should not generate harmful emissions or pollutants. & 0 & 0 & 0 & 1 & 1 \\
7.16. The system should have clear instructions for safe operation. & 0 & 1 & 0 & 3 & 4 \\
7.17. The system should be designed to minimize the risk of fire or hazards. & 0 & 0 & 0 & 0 & 0 \\
7.18. The system should ensure that all components are properly secured. & 0 & 1 & 0 & 1 & 2 \\
7.19. The system should be resistant to vandalism or accidental damage. & 0 & 0 & 0 & 1 & 1 \\
7.20. The system should have emergency shut-off procedures clearly documented. & 1 & 1 & 0 & 2 & 4 \\
\hline
\end{tabular}
\end{table}

% Table 2
\begin{table}[H]
\centering
\caption*{Table 2: CTQ needs prioritization.}
\label{table:2}
\rowcolors{2}{lightgray}{}
\begin{tabular}{|p{11cm}|c|c|c|c|c|}
\hline
\textbf{Need} & \textbf{\#1} & \textbf{\#2} & \textbf{\#3} & \textbf{\#4} & \textbf{Sum} \\
\hline
1.1. The system should minimize the time required to transfer water from the tank to the garden. & 3 & 3 & 4 & 3 & \cellcolor{yellow}13 \\
1.3. Water distribution should cover the entire garden area efficiently. & 2 & 2 & 3 & 2 & \cellcolor{yellow}9 \\
1.6. Water output should be consistent and reliable during the watering process. & 2 & 2 & 1 & 1 & 6 \\
1.12. The system should be easy to activate and deactivate when needed. & 0 & 1 & 1 & 1 & 3 \\
1.15. The system should minimize the physical effort required to operate it. & 1 & 1 & 0 & 1 & 3 \\
\hline
2.1. The system should provide sufficient pressure for effective water delivery. & 2 & 2 & 4 & 2 & \cellcolor{yellow}10 \\
2.10. The system should provide adequate flow rate for all watering methods. & 1 & 2 & 3 & 2 & \cellcolor{yellow}8 \\
\hline
5.10. The system should be designed to minimize noise pollution. & 1 & 1 & 4 & 1 & \cellcolor{yellow}7 \\
\hline
7.1. The system should comply with all safety standards. & 2 & 1 & 3 & 1 & \cellcolor{yellow}7 \\
\hline
\end{tabular}
\end{table}

% Table 3
\begin{table}[H]
\centering
\caption*{Table 3: Ranking of Critical to Quality (CTQ) needs.}
\label{table:3}
\rowcolors{2}{lightgray}{}
\begin{tabular}{|p{11cm}|c|c|c|c|c|}
\hline
\textbf{Need} & \textbf{\#1} & \textbf{\#2} & \textbf{\#3} & \textbf{\#4} & \textbf{Sum} \\
\hline
1.1. The system should minimize the time required to transfer water from the tank to the garden. & 6 & 6 & 6 & 6 & 24 \\
1.3. Water distribution should cover the entire garden area efficiently. & 3 & 4 & 2 & 5 & 14 \\
2.1. The system should provide sufficient pressure for effective water delivery. & 4 & 5 & 5 & 3 & 17 \\
2.10. The system should provide adequate flow rate for all intended watering methods. & 5 & 3 & 4 & 4 & 16 \\
5.10. The system should be designed to minimize noise pollution. & 2 & 2 & 3 & 2 & 9 \\
7.1. The system should comply with all relevant safety standards. & 1 & 1 & 1 & 1 & 4 \\
\hline
\end{tabular}
\end{table}

% Table 4
\begin{table}[H]
\centering
\caption*{Table 4: Final ranking of Critical to Quality (CTQ) needs.}
\label{table:4}
\rowcolors{2}{white}{white}
\begin{tabular}{|p{12cm}|c|}
\hline
\textbf{Need} & \textbf{Rank} \\
\hline
1.1. The system should minimize the time required to transfer water from the tank to the garden. & 1 \\
2.1. The system should provide sufficient pressure for effective water delivery. & 2 \\
2.10. The system should provide adequate flow rate for all watering methods. & 3 \\
1.3. Water distribution should cover the entire garden area efficiently. & 4 \\
5.10. The system should be designed to minimize noise pollution. & 5 \\
7.1. The system should comply with all safety standards. & 6 \\
\hline
\end{tabular}
\end{table}

\end{document}
